\documentclass[11pt]{article}
\input{preamble.tex}
\input{preamble-cours.tex}

% Solutions en fin de chapitre : \solutionsfalse pour les masquer.
\solutionstrue

\begin{document}

\chapitre{4}{Géométrie}

Dans tout le chapitre, les points sont dans le plan ou dans l'espace. $AB$ ou $|AB|$ désigne la longueur du segment $[AB]$. Un triangle $ABC$ est toujours non aplati : $A$, $B$, $C$ ne sont pas alignés.

\section{Notions utiles}\label{sec:notions}

\begin{definition}[Vecteur]\label{def:vecteur}
Un vecteur non nul est défini par une direction, un sens et une norme (sa longueur). Si $A\neq B$, on note $\vect{AB}$ le vecteur :
    \begin{itemize}
        \item de direction $(AB)$,
        \item de sens allant du point $A$ au point $B$,
        \item de longueur $AB$.
    \end{itemize}
Le vecteur $\vect{AA}$ est le vecteur nul $\vec{0}$. La norme d'un vecteur $\vec{u}$ se note $\|\vec{u}\|$ ou $|\vec{u}|$ ; ainsi $\|\vect{AB}\|=AB$.
Relation de Chasles : $\vect{AB} + \vect{BC} = \vect{AC}$. Si l'on fixe une origine $O$, on identifie un point $X$ au vecteur $\vect{OX}$ et on le note simplement $X$ : alors $\vect{AB} = B - A$.
\end{definition}

\begin{definition}[Produit scalaire]\label{def:produit-scalaire}
si $\vec{u}$ et $\vec{v}$ sont non nuls, $\vec{u}\cdot\vec{v}=\|\vec{u}\|\,\|\vec{v}\|\cos\theta$, où $\theta$ est l'angle entre $\vec{u}$ et $\vec{v}$ ; si l'un d'eux est nul, $\vec{u}\cdot\vec{v}=0$. On a $\vec{u}\cdot\vec{u}=\|\vec{u}\|^2$, et $\vec{u}\cdot\vec{v}=0$ si et seulement si $\vec{u}$ et $\vec{v}$ sont orthogonaux (ou l'un d'eux est nul). Le produit scalaire est symétrique ($\vec{u}\cdot\vec{v}=\vec{v}\cdot\vec{u}$) et se développe comme un produit de nombres, d'où $\|\vec{u}\pm\vec{v}\|^2=\|\vec{u}\|^2\pm2\,\vec{u}\cdot\vec{v}+\|\vec{v}\|^2$.
\end{definition}

\begin{definition}[Milieu d'un segment]\label{def:milieu}
Le milieu du segment $[AB]$ est le point $M$ tel que $\vect{AM} = \vect{MB}$. Dans un repère, ses coordonnées sont $ \left( \dfrac{x_A+x_B}{2} ; \dfrac{y_A+y_B}{2} ; \ldots \right)$ ; avec l'identification ci-dessus, $M=\dfrac{A+B}{2}$.
\end{definition}

\begin{definition}[Centroïde (centre de gravité)]\label{def:centroide}
Le centroïde de $n$ points $A_1, A_2, \ldots, A_n$ est l'unique point $S$ tel que
\[
\sum_{i=1}^{n} \vect{SA_i} = \vec{0},
\quad\text{c'est-à-dire}\quad S=\frac{A_1+\cdots+A_n}{n}.
\]
Pour un triangle, c'est le point d'intersection des médianes du triangle (les droites joignant un sommet au milieu du côté opposé).
\end{definition}

\begin{definition}[Parallélogramme]\label{def:parallelogramme}
Un quadrilatère $ABCD$, dont les quatre sommets ne sont pas alignés, est un parallélogramme si ses côtés opposés sont parallèles. Cela équivaut à $\vect{AB}=\vect{DC}$, et aussi au fait que les diagonales $[AC]$ et $[BD]$ ont le même milieu, c'est-à-dire $A+C=B+D$. Quatre points alignés qui vérifient $A+C=B+D$ forment un parallélogramme aplati.
\end{definition}

\begin{definition}[Médianes d'un quadrilatère]\label{def:medianes-quadrilatere}
dans un quadrilatère $ABCD$, ce sont les deux segments qui joignent les milieux de côtés opposés : celui de $[AB]$ à celui de $[CD]$, et celui de $[BC]$ à celui de $[DA]$. (À ne pas confondre avec les médianes d'un triangle.)
\end{definition}

\begin{definition}[Orthocentre et cercle circonscrit]\label{def:orthocentre}
l'orthocentre $H$ d'un triangle $ABC$ est le point d'intersection de ses trois hauteurs ; le centre $O$ de son cercle circonscrit est le point d'intersection des médiatrices de ses côtés, donc $OA=OB=OC$.
\end{definition}

\begin{proposition}[Orthocentre]\label{prop:orthocentre}
Si $O$ est pris pour origine, $H=A+B+C$.
\end{proposition}

\begin{proof}
En effet, le point $H'=A+B+C$ vérifie $\vect{AH'}\cdot\vect{BC}=(B+C)\cdot(C-B)=OC^2-OB^2=0$, donc $H'$ est sur la hauteur issue de $A$ ; de même, il est sur les deux autres hauteurs, donc $H'=H$.
\end{proof}

\begin{corollaire}[Droite d'Euler]\label{cor:euler}
Les points $O$, $G$, $H$ sont alignés et $\vect{OH}=3\,\vect{OG}$.
\end{corollaire}

\begin{proof}
Comme le centre de gravité est $G=\dfrac{A+B+C}{3}$, on obtient $\vect{OH}=3\,\vect{OG}$ : les points $O$, $G$, $H$ sont alignés (droite d'Euler).
\end{proof}

\begin{definition}[Angles]\label{def:angles}
$\angle BAC$ (ou $\widehat{BAC}$) désigne l'angle non orienté de sommet $A$ entre les demi-droites $[AB)$ et $[AC)$ ; il est compris entre $0^\circ$ et $180^\circ$. Un triangle est acutangle quand ses trois angles sont aigus (inférieurs à $90^\circ$).
\end{definition}

\begin{theoreme}[Angle inscrit]\label{thm:angle-inscrit}
soient $A$, $B$, $M$ trois points distincts d'un cercle de centre $O$. L'angle inscrit $\angle AMB$ vaut la moitié de l'angle au centre qui intercepte le même arc $AB$ (l'arc qui ne contient pas $M$ ; si c'est le grand arc, cet angle au centre dépasse $180^\circ$). Conséquences : deux points $M$ et $N$ du cercle situés du même côté de la droite $(AB)$ vérifient $\angle AMB=\angle ANB$ ; et, réciproquement, si deux points $M$ et $N$ sont du même côté de $(AB)$ et vérifient $\angle AMB=\angle ANB$, alors $A$, $B$, $M$, $N$ sont sur un même cercle.
\end{theoreme}

\begin{definition}[Triangles semblables]\label{def:semblables}
deux triangles sont semblables s'ils ont les mêmes angles, ce qui équivaut à avoir des côtés proportionnels. Il suffit aussi qu'ils aient un angle égal compris entre deux côtés proportionnels. Deux triangles semblables $XYZ$ et $X'Y'Z'$ (sommets dans cet ordre) sont de même sens si l'on tourne dans le même sens en parcourant $X\to Y\to Z$ et $X'\to Y'\to Z'$, et de sens opposés sinon.
\end{definition}

\section{À retenir}\label{sec:retenir}

\begin{retenir}
\begin{itemize}
  \item $\vect{AB}=B-A$ et, avec une origine fixée, un point s'identifie au vecteur qui le joint à l'origine ; relation de Chasles $\vect{AB}+\vect{BC}=\vect{AC}$.
  \item $\vec{u}\cdot\vec{u}=\|\vec{u}\|^2$ ; $\vec{u}\cdot\vec{v}=0$ si et seulement si $\vec{u}$ et $\vec{v}$ sont orthogonaux ; $\|\vec{u}\pm\vec{v}\|^2=\|\vec{u}\|^2\pm2\,\vec{u}\cdot\vec{v}+\|\vec{v}\|^2$.
  \item Milieu de $[AB]$ : $\dfrac{A+B}{2}$ ; centroïde de $A_1,\ldots,A_n$ : $\dfrac{A_1+\cdots+A_n}{n}$.
  \item $ABCD$ est un parallélogramme si et seulement si $\vect{AB}=\vect{DC}$, c'est-à-dire $A+C=B+D$.
  \item Si $O$, centre du cercle circonscrit, est l'origine, l'orthocentre est $H=A+B+C$ ; donc $\vect{OH}=3\,\vect{OG}$ (droite d'Euler).
  \item Angle inscrit : moitié de l'angle au centre ; deux points d'un même côté de $(AB)$ voient $[AB]$ sous le même angle si et seulement si ils sont sur un même cercle avec $A$ et $B$.
  \item Deux triangles sont semblables s'ils ont les mêmes angles, ou des côtés proportionnels, ou un angle égal compris entre deux côtés proportionnels.
\end{itemize}
\end{retenir}

\section{Exercices}\label{sec:exercices}

\exochap{recueil-4-01-milieux-quadrilatere}{Les milieux d'un quadrilatère}

\exochap{recueil-4-02-hexagone-centroides}{Hexagone formé par les centroïdes}

\exochap{recueil-4-03-centroide-n-points}{Centroïde de $n$ points}

\exochap{recueil-4-04-diagonales-orthogonales}{Diagonales orthogonales}

\exochap{recueil-4-05-relation-espace}{Relation vectorielle dans l'espace}

\exochap{recueil-4-06-bhoc-parallelogramme}{Orthocentre, centre circonscrit et angle de $60^\circ$}

\exochap{recueil-4-07-medianes-relation}{Somme des carrés et centre de gravité}

\exochap{recueil-4-08-triangle-isocele}{Bissectrice d'un triangle isocèle}

\section{Problèmes}\label{sec:problemes}

\probleme{recueil-4-p1-relation-diagonales}{Relation sur les diagonales}

\probleme{recueil-4-p2-rectangle-espace}{Rectangle dans l'espace}

\probleme{recueil-4-p3-theoreme-euler}{Théorème d'Euler}

\probleme{recueil-4-p4-produit-scalaire}{Produit scalaire dans un rectangle}

\probleme{recueil-4-p5-somme-vecteurs}{Somme de vecteurs bornés}

\probleme{recueil-4-p6-triangles-equilateraux}{Triangles équilatéraux et centre de gravité}

\probleme{recueil-4-p7-hexagone-inscrit}{Hauteur, bissectrice et rayon circonscrit}

\probleme{recueil-4-p8-point-fixe}{Une tangente au cercle $AMC$}

\probleme{recueil-4-p9-milieux}{Deux triangles semblables}

\end{document}
